\section{Related literature}\label{sec:lit}

\para{Analyst revisions and target prices:} Earnings-forecast revisions predict returns \citep{stickel1991}. Revision-based momentum and price momentum are related but different \citep{chan1996}. The market under-reacts to revisions that are less visible \citep{gleason2003}, and recommendation changes carry information that overlaps with momentum \citep{jegadeesh2004}. Target prices move prices \citep{brav2003,asquith2005}. Analysts' ability to forecast the level of a target price is limited and does not persist \citep{bradshaw2013}, and relative valuation inside an industry is more informative than absolute upside \citep{da2011}. The same pattern is found here: signals on target-price levels fail, signals on the flow of revisions work.

\para{Leaders, herders and disagreement:} \citet{cooper2001} identify lead analysts from the timing of their forecasts; lead revisions have more price impact. In herding models and evidence \citep{welch2000,hong2000,clement2005}, revisions issued close to other revisions carry less information. LMH makes this timing distinction into a stock-level cross-sectional signal. \citet{diether2002} link dispersion of opinion to returns, which gives one of the comparables.

\para{Institutional factor strategies:} The comparables are the standard toolkit of cross-sectional momentum \citep{jegadeesh1993}, short-term reversal, low volatility, betting-against-beta \citep{frazzini2014}, residual momentum \citep{blitz2011}, proximity to the 52-week high \citep{george2004} and time-series momentum \citep{moskowitz2012}. Volatility management \citep{barroso2015,moreira2017} and regime filters \citep{hamilton1989} are the usual overlays.

\para{Construction, costs and inference:} Overlapping holding-period portfolios date from \citet{jegadeesh1993}. Many anomalies do not survive trading costs \citep{novymarx2016}, and estimation error favours naive weights \citep{demiguel2009}. \citet{harvey2016} argue for higher hurdles in inference. The tools used here are the probabilistic and deflated Sharpe ratios \citep{bailey2012,bailey2014}, the probability of backtest overfitting \citep{bailey2017}, Holm's correction \citep{holm1979}, the stationary bootstrap \citep{politis1994}, Sharpe-difference tests \citep{jobson1981,memmel2003}, Newey--West and Fama--MacBeth inference \citep{newey1987,fama1973}, asymmetric GARCH \citep{glosten1993} and purged evaluation \citep{lopezdeprado2018}.

\para{Data snooping, small samples and generative scenarios:} \citet{white2000}, \citet{hansen2005} and \citet{romano2005} give bootstrap tests that allow for a search over many strategies; \citet{ledoit2008} give a robust test for differences of Sharpe ratios. Multiverse analysis \citep{steegen2016} and the specification curve \citep{simonsohn2020} report results over all reasonable analytic choices. \citet{lo2002} corrects Sharpe ratios for serial correlation, and \citet{lo2008} separates the active and passive parts of a strategy's return. On the Bayesian side I use the Bayesian bootstrap \citep{rubin1981}, hierarchical models with weakly informative priors on variance components \citep{gelman2006} and empirical-Bayes shrinkage \citep{efron2010}; \citet{mclean2016} motivate sceptical priors with the post-publication decay of anomalies. For path simulation the industry standard is filtered historical simulation \citep{baroneadesi1999} with variance targeting \citep{engle1996}. Deep generative models for financial time series include GANs \citep{wiese2020} and TimeGAN \citep{yoon2019}. Denoising diffusion models \citep{ho2020,nichol2021,song2021} lead generative modelling in other fields. The approach here combines the econometric and diffusion approaches in a form that works with small samples.
